% ============================================================
% CAPITOLO 3
% Data Governance e intelligenza artificiale negli atenei europei:
% framework, architetture e decisioni strategiche
% ============================================================
\chapter{Data Governance e intelligenza artificiale negli atenei europei}
\label{chap:data-governance}

Le istituzioni universitarie europee occupano una posizione peculiare nel panorama della
Pubblica Amministrazione: sono al contempo produttori e custodi di dati ad alta densità
strategica, operatori soggetti a obblighi normativi stratificati e, sempre più frequentemente,
sperimentatori di soluzioni di intelligenza artificiale applicate a processi decisionali istituzionali.
Questa triplice condizione --- custode del dato, soggetto regolato, innovatore --- rende gli atenei un
laboratorio privilegiato per osservare le tensioni tra le esigenze di sovranità digitale, i vincoli
di compliance e i modelli organizzativi richiesti dall'IA in produzione.

Il presente capitolo analizza tale posizione a partire da un'angolatura manageriale. Nella prima
sezione si ricostruisce la specificità degli atenei come organizzazioni \textit{data-intensive} nella PA.
Nella seconda si esaminano i principali \textit{framework} internazionali di \textit{data governance} ---
DAMA-DMBOK, ISO/IEC\,38505 e ISO/IEC\,27001 --- nella loro applicabilità al contesto accademico.
La terza sezione introduce i \textit{maturity model} come strumenti decisionali e propone un modello
sintetico di \textit{AI-readiness} per gli atenei. La quarta sezione analizza tre casi europei
verificati di IA \textit{on-premise} o ibrida in ambito accademico. La quinta sezione distilla le
implicazioni strategiche e traccia il raccordo verso l'analisi comparativa del Capitolo\,4.

% ============================================================
\section{Gli atenei come organizzazioni \textit{data-intensive} nella Pubblica Amministrazione}
\label{sec:atenei-data-intensive}

\subsection{La stratificazione del patrimonio informativo universitario}

Nell'ambito della PA digitale, le università presentano una struttura informativa che non trova
riscontro immediato in altre categorie di enti pubblici. Il dato universitario si articola su tre
livelli sovrapposti, ciascuno con proprie esigenze di protezione, valorizzazione e governance.

Il primo livello --- quello dei \textbf{dati operativi di massa} --- comprende le carriere degli
studenti, i pagamenti, le presenze, le iscrizioni e i dati anagrafici. Si tratta di informazioni
ad alto volume, trattate attraverso piattaforme centralizzate come ESSE3, U-GOV e i sistemi CINECA,
che le università italiane adottano in forma consortile. Il secondo livello --- quello dei
\textbf{dati sensibili} --- include le condizioni economiche degli studenti (rilevanti per
l'attribuzione di borse di studio), le disabilità, i dati sanitari connessi a sperimentazioni
cliniche e, in taluni atenei, i dati biometrici raccolti nell'ambito di trial di ricerca. Il
terzo livello --- quello dei \textbf{dati di ricerca ad alto valore strategico} --- comprende i
\textit{dataset} proprietari, i brevetti in iter di registrazione, i risultati preliminari di
ricerca finanziata da Horizon Europe e, sempre più, i modelli di \textit{machine learning}
addestrati su dati istituzionali \parencite{McNicol2024DataEthics}.

Questa stratificazione non è neutra rispetto alle scelte architetturali. Come osservano
\textcite{Aguilar2026}, la \textit{data governance} centrata sul dato --- anziché limitata alla
sola sicurezza informatica --- costituisce la precondizione per la creazione di valore
organizzativo responsabile nelle PA europee. Applicato al contesto universitario, questo principio
implica che la scelta di un modello di \textit{deployment} dell'IA non possa prescindere dalla
natura e dalla classificazione del dato trattato: un sistema di supporto alla valutazione degli
studenti opera su dati del secondo livello; un modello predittivo per la ricerca può coinvolgere
dati del terzo. In entrambi i casi, la dipendenza da infrastrutture \textit{cloud} esterne espone
l'ateneo a rischi di perdita di controllo sul patrimonio informativo che la normativa vigente ---
in primo luogo il GDPR e l'AI Act --- non consente di sottovalutare.

\subsection{Obblighi normativi specifici degli atenei: AI Act Annex III e GDPR}

Il quadro regolatorio europeo ha introdotto, con l'AI Act, un elemento di discontinuità
rilevante per gli atenei: la classificazione come \textbf{sistemi ad alto rischio} dei sistemi
di IA utilizzati in ambito educativo e formativo. L'Allegato III del Regolamento (UE) 2024/1689
\parencite{AIAct2024}, al punto\,3, elenca tra i sistemi \textit{high-risk} quelli impiegati per
«determinare l'accesso o l'ammissione a istituti di istruzione e formazione professionale»,
per «valutare i partecipanti a test e fissare risultati di apprendimento» e per «rilevare
comportamenti vietati durante esami». Ne consegue che qualunque ateneo europeo che adotti sistemi
di IA per l'ammissione ai corsi di laurea magistrale, per il riconoscimento automatico dei crediti
o per la valutazione adattiva diventa \textit{deployer} di un sistema \textit{high-risk} e come
tale è soggetto agli obblighi di registrazione nel database UE, di valutazione d'impatto, di
supervisione umana e di conservazione dei \textit{log} per almeno dieci anni
\parencite{Hohmann2025AIActGDPR}.

Il punto cruciale, da un punto di vista manageriale, è che tali obblighi richiedono un
\textbf{controllo diretto sull'infrastruttura}. La supervisione umana non è verificabile se i
dati di inferenza transitano per server di fornitori terzi privi di obblighi di trasparenza
equivalenti. La conservazione dei \textit{log} non è garantibile se l'architettura \textit{cloud}
del fornitore non prevede la portabilità e l'esportazione dei \textit{log} stessi. L'on-premise
--- o l'architettura ibrida con componente di inferenza \textit{on-premise} --- è dunque una
risposta strutturale a un vincolo normativo, non una preferenza tecnologica. Come argomentano
\textcite{Hohmann2025AIActGDPR}, l'AI Act costituisce uno strato regolatorio complementare e più
stringente rispetto al GDPR, che insieme impongono una governance del dato algoritmica che è
più agevolmente implementata in ambienti a pieno controllo istituzionale.

\subsection{La dimensione strategica: università come \textit{data steward} istituzionali}

Al di là degli obblighi normativi, esiste una ragione strategica per cui gli atenei europei
stanno ripensando il proprio rapporto con l'infrastruttura dati. La ricerca di \textcite{NunezCanal2026}
dimostra empiricamente che la competenza del corpo docente in materia di IA --- intesa come
capacità di gestire dati, valutare modelli e applicare principi di etica computazionale --- è
correlata positivamente alla capacità istituzionale dell'ateneo di innovare i propri processi
decisionali. Questo risultato suggerisce che la \textbf{governance del dato} non è una funzione
tecnica separata dalla strategia accademica: è una capacità organizzativa abilitante
(\textit{organizational capability}) che determina la qualità delle decisioni istituzionali
\parencite{Dabbous2026AIReadiness}.

Il concetto di \textit{data stewardship} --- la responsabilità istituzionale di gestire i dati
come un patrimonio pubblico, con obblighi di accuratezza, accessibilità e protezione --- è
codificato nel DAMA-DMBOK come una delle undici aree di conoscenza della \textit{data governance}
\parencite{DAMA_DMBOK}. Applicato agli atenei, il \textit{data stewardship} implica che le
scelte infrastrutturali siano governate da un organo con mandato istituzionale --- tipicamente
il \textbf{Delegato alla Transizione Digitale} previsto dall'art.\,17 del CAD --- e non
delegate de facto al fornitore tecnologico. L'ateneo che esternalizza l'infrastruttura IA senza
un'adeguata contrattualizzazione dei diritti sui dati e dei \textit{log} di inferenza non sta
compiendo una scelta tecnologica: sta cedendo la propria capacità di \textit{data stewardship},
con conseguenze dirette sulla compliance normativa e sulla sovranità istituzionale.

% ============================================================
\section{Framework internazionali di \textit{data governance}: DAMA, ISO/IEC\,38505, ISO/IEC\,27001}
\label{sec:framework-governance}

\subsection{DAMA-DMBOK: la mappa delle responsabilità organizzative}

Il \textit{Data Management Body of Knowledge} (DAMA-DMBOK), nella sua seconda edizione del 2017
\parencite{DAMA_DMBOK}, rappresenta il riferimento più diffuso per la strutturazione della
\textit{data governance} a livello internazionale. Il \textit{framework} organizza le funzioni
di gestione del dato attorno a undici aree di conoscenza, tra le quali le più rilevanti per
il contesto universitario \textit{on-premise} sono: l'\textbf{architettura dei dati}
(\textit{Data Architecture}), che definisce dove risiedono i dati e come si relazionano tra
loro; la \textbf{sicurezza dei dati} (\textit{Data Security}), che governa accessi, cifratura
e tracciabilità; il \textbf{management dei metadati} (\textit{Metadata Management}), che
garantisce la tracciabilità delle trasformazioni subite dai dati --- incluse quelle operate da
algoritmi di IA.

Quest'ultima area assume un rilievo particolare alla luce dell'AI Act: la documentazione tecnica
obbligatoria per i sistemi \textit{high-risk} (art.\,11) richiede esattamente la tracciabilità
dei \textit{dataset} utilizzati nell'addestramento, che il \textit{Metadata Management} del
DAMA-DMBOK è progettato per garantire. Un ateneo che non abbia implementato tale area --- come
è il caso della maggioranza degli atenei italiani, stando alle analisi di maturità disponibili
--- non è in grado di documentare in modo conforme la catena di provenienza dei dati utilizzati
per addestrare modelli applicati alla valutazione degli studenti
\parencite{McNicol2024DataEthics}.

Dalla prospettiva della governance organizzativa, il DAMA-DMBOK assegna le responsabilità
secondo una struttura a tre livelli: il livello \textbf{strategico} (CDA, Senato Accademico,
Rettore) definisce le politiche; il livello \textbf{tattico} (CIO, Delegato alla Transizione
Digitale, DPO) traduce le politiche in processi; il livello \textbf{operativo} (amministratori
di sistema, data steward di dipartimento) implementa i processi e segnala le anomalie.
La coerenza tra questi tre livelli è la condizione necessaria per la sostenibilità di un
sistema di IA \textit{on-premise}, che richiede competenze distribuite e non centralizzate
in un singolo fornitore esterno.

\subsection{ISO/IEC\,38505: la governance del dato come funzione di \textit{board}}

La norma ISO/IEC\,38505-1:2017 \parencite{ISO38505} estende al dominio della \textit{data
governance} il modello \textit{Evaluate--Direct--Monitor} già consolidato dall'IT Governance
ISO/IEC\,38500. Il contributo specifico di questa norma rispetto al DAMA-DMBOK è il suo
orientamento esplicito al \textbf{livello di \textit{board}}: la governance del dato non è
una funzione tecnica ma una responsabilità di vertice, analoga alla governance finanziaria.

Applicato agli atenei, questo principio implica che il Consiglio di Amministrazione valuti
(\textit{Evaluate}) le politiche di utilizzo del dato istituzionale --- incluse le decisioni
sul modello di \textit{deployment} dell'IA ---; che indirizzi (\textit{Direct}) il management
nella loro attuazione; e che monitorizzi (\textit{Monitor}) i risultati attraverso indicatori
definiti. Il Nucleo di Valutazione, previsto dalla normativa ANVUR, può svolgere la funzione
di \textit{Monitor} esterno, verificando la coerenza tra le dichiarazioni dell'ateneo in
materia di \textit{data governance} e le pratiche effettive \parencite{ISO38505}.

La norma ISO/IEC\,38505 è particolarmente rilevante perché fornisce il fondamento per
trattare la scelta \textit{on-premise} come una \textbf{decisione di governance}, non come
una scelta tecnica. Il CDA che delibera di adottare un'infrastruttura \textit{on-premise}
per i sistemi di IA che trattano dati sensibili degli studenti sta esercitando la funzione
di \textit{Evaluate} prevista dalla norma, assumendosi la responsabilità istituzionale
che l'AI Act attribuisce al \textit{deployer}.

\subsection{ISO/IEC\,27001:2022 e la sicurezza come prerequisito dell'IA}

La norma ISO/IEC\,27001, nella sua edizione 2022 \parencite{ISO27001}, introduce rilevanti
novità rispetto alla versione 2013, consolidando 114 controlli in 93 e aggiungendo nuovi
domini relativi alla sicurezza \textit{cloud}, alla \textit{threat intelligence} e alla
gestione della continuità operativa in ambienti ibridi. Per gli atenei che trattano dati
soggetti alla Direttiva NIS\,2 --- tra cui figurano, in base al perimetro italiano, le
università che gestiscono infrastrutture di ricerca critiche o che partecipano al Polo
Strategico Nazionale --- la certificazione ISO/IEC\,27001 è diventata un requisito di fatto.

Nel contesto dell'IA \textit{on-premise}, la norma assume un ruolo di
\textbf{prerequisito infrastrutturale}: un sistema di IA che inferisce su dati personali in
un ambiente non certificato secondo ISO/IEC\,27001 non soddisfa i requisiti di sicurezza
dell'AI Act per i sistemi \textit{high-risk}. Viceversa, un ateneo certificato dispone già
del sistema di gestione della sicurezza delle informazioni (ISMS) necessario per documentare
e auditare i processi di \textit{training}, \textit{fine-tuning} e inferenza dei modelli.

Come dimostrano \textcite{PatonRomero2018GreenIT} nell'analisi comparativa dei \textit{maturity
model} basati su standard ISO, l'adozione progressiva degli standard internazionali di
governance IT segue una logica cumulativa: ogni livello di maturità raggiunto costituisce la
base per il livello successivo. Questa logica si applica direttamente agli atenei: la
certificazione ISO/IEC\,27001 non è un adempimento isolato ma il tassello di un percorso
verso l'\textit{AI-readiness} istituzionale.

La tabella seguente sintetizza il confronto tra i tre \textit{framework} analizzati.

\begin{table}[H]
\centering
\caption{Confronto tra i principali \textit{framework} di \textit{data governance} per gli atenei.}
\label{tab:framework-governance}
\begin{tabularx}{\textwidth}{l X X X X}
\toprule
\textbf{Framework} & \textbf{Livello primario} & \textbf{Applicabilità PA} & \textbf{Integrazione AI Act} & \textbf{Contributo on-premise} \\
\midrule
DAMA-DMBOK v2 & Operativo-tattico & Alta (data management) & Metadata Management per Annex III & Inventario e classificazione del dato \\
ISO/IEC 38505 & Strategico (board) & Alta (governance pubblica) & Framework Evaluate-Direct-Monitor & Legittimazione della scelta architetturale \\
ISO/IEC 27001 & Tattico-operativo & Obbligatoria (NIS2) & Prerequisito ISMS per high-risk AI & Certificazione dell'infrastruttura di hosting \\
\bottomrule
\end{tabularx}
\end{table}

% ============================================================
\section{\textit{Maturity model} per la \textit{AI-readiness} universitaria}
\label{sec:maturity-model}

\subsection{La maturità digitale come variabile manageriale}

Il concetto di \textbf{maturità digitale} --- inteso come la capacità di un'organizzazione di
adottare, integrare e valorizzare le tecnologie digitali in modo sistematico e misurabile ---
ha assunto negli ultimi anni un ruolo centrale nella letteratura di \textit{Information Systems}
e nel dibattito sulla trasformazione della PA europea. I \textit{maturity model} non sono mere
classificazioni descrittive: sono \textbf{strumenti decisionali} che consentono al management
di identificare il gap tra lo stato corrente e lo stato obiettivo, di pianificare interventi
prioritari e di allocare risorse in modo coerente con la traiettoria di sviluppo istituzionale
\parencite{Torichnyi2025StrategicMgmt}.

Nel caso specifico degli atenei, la maturità digitale condiziona direttamente la fattibilità
dell'IA \textit{on-premise}: un ateneo privo di un inventario strutturato del proprio patrimonio
dati, di un sistema di gestione della sicurezza certificato e di competenze interne di
\textit{data engineering} non è in grado di operare e mantenere un sistema di IA su
infrastruttura propria. Il \textit{maturity model} consente di rendere questo gap visibile e
quantificabile, trasformandolo da ostacolo percepito a piano di sviluppo programmabile.

\subsection{Modelli di riferimento esistenti}

Il \textbf{Data Management Maturity Model (DMM)} del CMMI Institute \parencite{CMMI}
rappresenta il riferimento più specifico per la maturità della gestione del dato. Articolato
in cinque livelli di maturità e sei categorie di processo --- \textit{Data Management Strategy,
Data Governance, Data Quality, Platform \& Architecture, Data Operations, Supporting Processes}
--- il DMM consente di diagnosticare la capacità organizzativa di un ateneo su ciascuna
dimensione in modo indipendente. Un ateneo può essere al livello\,3 (\textit{Defined}) per
la \textit{Data Quality} e al livello\,1 (\textit{Performed}) per la \textit{Platform \&
Architecture}: questa asimmetria è informativa per il management, perché indica dove concentrare
gli investimenti.

Il modello della \textbf{maturità della \textit{records governance}} proposto da
\textcite{Abdallah2026ARMA}, basato sul framework ARMA e applicato in un contesto di PA locale,
evidenzia che la maturità documentale è un predittore affidabile della maturità digitale
complessiva: le organizzazioni che non governano efficacemente i propri documenti strutturati
difficilmente governano i \textit{dataset} necessari all'IA. Questo risultato, pur proveniente
da un contesto geografico diverso, si generalizza al caso degli atenei italiani, molti dei
quali presentano sistemi di \textit{records management} frammentati tra dipartimenti e
infrastrutture ereditate.

L'\textbf{Open Data Maturity Index} della Commissione Europea, nella sua edizione 2024
\parencite{DESI2024}, misura il livello di maturità nella pubblicazione di dati aperti degli
Stati membri e delle loro PA. L'Italia si colloca nel cluster \textit{Fast Tracker} per la
terza edizione consecutiva, con un punteggio complessivo superiore alla media UE, ma con
marcate asimmetrie tra le PA centrali --- che guidano la classifica --- e le PA locali e gli
atenei, che presentano livelli di maturità significativamente inferiori nella dimensione
della \textit{data quality} e del \textit{data reuse}.

\subsection{Un modello sintetico di \textit{AI-readiness} per gli atenei: proposta dell'autore}

A partire dai modelli DMM, DAMA-DMBOK e ISO/IEC\,38505, si propone un \textbf{modello sintetico
di \textit{AI-readiness}} a quattro livelli, calibrato sulle specificità degli atenei europei
che intendano adottare sistemi di IA con componente \textit{on-premise}. Il modello non
sostituisce i \textit{framework} di riferimento ma li sintetizza in una griglia diagnostica
applicabile con risorse limitate, orientata al decision-making del management accademico.

\begin{table}[H]
\centering
\caption{Modello sintetico di \textit{AI-readiness} per gli atenei europei (proposta dell'autore).}
\label{tab:ai-readiness}
\begin{tabularx}{\textwidth}{l l X X}
\toprule
\textbf{Livello} & \textbf{Label} & \textbf{Caratteristiche organizzative} & \textbf{Modello AI compatibile} \\
\midrule
1 & \textit{Ad hoc} & Nessuna governance strutturata del dato; assenza di DPO operativo; infrastruttura IT non inventariata & Nessun deployment autonomo; solo SaaS cloud \\
2 & \textit{Managed} & Politiche di sicurezza esistenti ma non integrate; DPO nominato; primo inventario dei sistemi informativi & Cloud qualificato; nessuna componente on-premise \\
3 & \textit{Defined} & DAMA o ISO 27001 adottati; processi documentati; Delegato Transizione Digitale attivo; dati classificati & Architettura ibrida possibile; inferenza on-premise su dati sensibili \\
4 & \textit{AI-Ready} & ISO 27001 certificata; DMM livello 3+; competenze MLOps interne; infrastruttura HPC consortile & On-premise sostenibile; federated learning inter-istituzionale \\
\bottomrule
\end{tabularx}
\end{table}

Il passaggio dal livello\,2 al livello\,3 è il salto organizzativo più critico: richiede non solo
investimenti tecnologici ma un cambio di cultura istituzionale, con il dato che cessa di essere
un sottoprodotto amministrativo e diventa un \textbf{asset strategico governato}. Il passaggio
dal livello\,3 al livello\,4 è prevalentemente tecnologico e finanziario: presuppone che il
cambio culturale e organizzativo sia già avvenuto e che l'ateneo disponga --- o abbia accesso
consortile --- a infrastrutture HPC adeguate all'addestramento e all'esecuzione di modelli
di IA in produzione.

% ============================================================
\section{Casi europei di IA \textit{on-premise} e ibrida negli atenei: analisi comparativa}
\label{sec:casi-europei}

\subsection{Metodologia di selezione}

I casi analizzati in questa sezione sono stati selezionati sulla base di tre criteri: (a)
\textbf{verificabilità delle fonti} --- sono stati inclusi solo atenei con documentazione
pubblica accessibile (pubblicazioni accademiche, report istituzionali, \textit{dataset} aperti);
(b) \textbf{rilevanza architetturale} --- i casi presentano componenti \textit{on-premise} o
ibride significative, non solo utilizzo di piattaforme SaaS; (c) \textbf{pertinenza manageriale}
--- ogni caso illustra almeno uno dei tre cardini della presente tesi: sovranità digitale,
compliance normativa o modello organizzativo.

\subsection{Caso A --- KU Leuven (Belgio): \textit{federated learning} per la ricerca biomedica e la
protezione del dato paziente}

La Katholieke Universiteit Leuven è uno degli atenei europei più attivi nello sviluppo di
architetture di \textit{federated learning} applicate alla ricerca biomedica. Il contesto
normativo è quello della ricerca clinica multi-istituzionale: i dati dei pazienti --- regolati
dal GDPR e, per la ricerca sperimentale, dai requisiti dei comitati etici nazionali --- non
possono essere centralizzati in un'unica infrastruttura senza il consenso esplicito degli
interessati, consenso che nella pratica della ricerca distribuita è spesso impraticabile
da ottenere.

La risposta architetturale adottata da KU Leuven, in collaborazione con partner di ricerca
europei nell'ambito di progetti Horizon Europe, è il \textit{federated learning}: i modelli
di IA vengono addestrati \textbf{localmente} su ciascun nodo ospedaliero o universitario,
e solo i \textit{gradient updates} --- aggregati e resi anonimi --- vengono condivisi con
il coordinatore del consorzio \parencite{GadekalluFL2026}. Il dato grezzo non abbandona mai
il perimetro fisico e giuridico dell'istituzione che lo possiede: è questa la garanzia di
sovranità che rende il modello conforme al GDPR senza richiedere il consenso esplicito per
ogni singolo dato.

Dal punto di vista manageriale, l'esperienza di KU Leuven illustra un principio generale: la
\textbf{sovranità del dato non è un vincolo alla collaborazione scientifica ma la condizione
che la rende possibile} a scala multi-istituzionale. Il \textit{federated learning} risolve
il paradosso tra la necessità di \textit{dataset} sufficientemente grandi per addestrare
modelli robusti e la proibizione di concentrare dati sensibili in un unico repository. Questa
logica è direttamente replicabile in ambito universitario non biomedico: sistemi di supporto
alla didattica che apprendono dai comportamenti degli studenti di più atenei, senza centralizzare
i dati di nessuno, sono tecnicamente e normativamente praticabili con le stesse architetture
\parencite{GadekalluFL2026}.

\subsection{Caso B --- TU Delft (Paesi Bassi): \textit{learning analytics} ibrida e GDPR by design}

Il Politecnico di Delft ha sviluppato, a partire dal 2019, un sistema di \textit{learning
analytics} per la predizione del rischio di abbandono (\textit{dropout prediction}) tra gli
studenti del primo anno di ingegneria. Il sistema integra dati accademici --- frequenza,
voti intermedi, partecipazione ai \textit{tutorial} --- con dati comportamentali provenienti
dalla piattaforma LMS (\textit{Learning Management System}). La sensibilità di questi dati
--- che rivelano pattern comportamentali individuali e possono essere utilizzati per inferire
condizioni di disagio psicologico o difficoltà economiche --- ha reso necessaria una scelta
architetturale orientata al \textit{privacy by design} \parencite{Wichmann2020GDPR}.

L'architettura adottata da TU Delft è \textbf{ibrida con componente di inferenza
\textit{on-premise}}: i dati grezzi degli studenti risiedono su server dell'ateneo certificati
ISO/IEC\,27001; il processo di \textit{feature engineering} e \textit{model training} avviene
su infrastruttura locale; solo le predizioni aggregate --- non i profili individuali --- vengono
rese disponibili ai tutor attraverso un'interfaccia \textit{cloud}. Questo disegno garantisce
che i dati personali non abbandonino mai il perimetro dell'ateneo, soddisfacendo il requisito
di \textit{data minimisation} del GDPR (art.\,5, par.\,1, lett.\,c) e il principio di
\textit{privacy by design} (art.\,25).

Il caso TU Delft è significativo perché dimostra che l'architettura ibrida non è solo una
soluzione per contesti ad alto rischio come la ricerca biomedica: è praticabile e sostenibile
anche per applicazioni didattiche ordinarie, purché il management accademico abbia chiarito
ex ante la classificazione dei dati trattati e le responsabilità del \textit{deployer}
\parencite{Hohmann2025AIActGDPR}. Il sistema di TU Delft, classificabile come \textit{high-risk}
ex AI Act Annex III (valutazione dei partecipanti a percorsi di formazione), richiede ora la
documentazione tecnica e la registrazione nel database UE previste dal Regolamento: l'architettura
\textit{on-premise} per la componente di \textit{training} facilita significativamente questo
adempimento.

\subsection{Caso C --- CINECA e gli atenei italiani: il modello di HPC consortile come
\textit{on-premise} collettivo}

Il Consorzio CINECA rappresenta il caso più rilevante per il contesto italiano e, per certi
aspetti, uno dei modelli più avanzati in Europa di \textit{on-premise} collettivo per la
ricerca universitaria. Fondato nel 1969 e partecipato da oltre settanta atenei italiani e da
enti pubblici di ricerca, CINECA gestisce il \textbf{supercomputer Leonardo}, installato
presso la sede di Casalecchio di Reno (Bologna) nel 2022 e classificato tra i dieci sistemi
HPC più potenti al mondo secondo la lista Top500 \parencite{EuroHPC_Leonardo}.

Il modello CINECA è, da un punto di vista architetturale, un'infrastruttura \textit{on-premise}
governata da un consorzio di istituzioni pubbliche italiane, con piena sovranità giuridica
nazionale sui dati e sui sistemi. Gli atenei aderenti accedono alle risorse computazionali
attraverso un sistema di allocazione a progetto, simile al modello dei \textit{cloud provider}
ma con la differenza fondamentale che l'infrastruttura fisica è italiana, il soggetto gestore è
un ente pubblico non profit e i dati trattati non sono soggetti alla giurisdizione di stati
extra-UE. Il \textbf{CLOUD Act} statunitense \parencite{CLOUDAct_Wire2025}, che obbliga i fornitori
\textit{cloud} con sede negli USA a consegnare i dati alle autorità americane su richiesta,
non si applica a CINECA.

Dal punto di vista manageriale, il modello CINECA offre agli atenei italiani medi ---
che non dispongono delle risorse per infrastrutture GPU proprietarie di scala ---
la possibilità di operare sistemi di IA in modalità \textit{on-premise} collettivo,
condividendo i costi fissi dell'infrastruttura e mantenendo la sovranità sul dato.
\textcite{Madhavan2025Storage} identificano nel modello consortile una delle configurazioni
ottimali per bilanciare costi finanziari, compliance normativa e prestazioni tecniche in
ambito AI, in quanto consente di sfruttare le economie di scala dell'HPC senza cedere il
controllo dell'infrastruttura a fornitori terzi. Il sistema Leonardo supporta già oggi
applicazioni di \textit{machine learning} in ambito climatico, farmacologico e di fisica
delle particelle per gli atenei del consorzio: l'estensione a modelli di IA per la
governance universitaria è tecnicamente immediata.

\subsection{Analisi comparativa: pattern ricorrenti e lezioni manageriali}

I tre casi analizzati mostrano configurazioni architetturali differenti --- \textit{federated}
(KU Leuven), ibrida (TU Delft), consortile (CINECA) --- ma convergono su alcuni pattern
ricorrenti che costituiscono le lezioni manageriali fondamentali del capitolo.

\begin{table}[H]
\centering
\caption{Analisi comparativa dei casi europei di IA \textit{on-premise} negli atenei.}
\label{tab:casi-europei}
\begin{tabularx}{\textwidth}{l X X X X}
\toprule
\textbf{Ateneo/Ente} & \textbf{Paese} & \textbf{Modello} & \textbf{Applicazione AI} & \textbf{AI-Readiness} \\
\midrule
KU Leuven & Belgio & Federated on-premise & Ricerca biomedica multi-istituzionale & Livello 4 \\
TU Delft & Paesi Bassi & Ibrido (inferenza on-premise) & Learning analytics, dropout prediction & Livello 3--4 \\
CINECA/IT & Italia & On-premise consortile & HPC per ricerca e governance & Livello 4 (infrastruttura) \\
\bottomrule
\end{tabularx}
\end{table}

Il primo pattern è la \textbf{normativa come driver architetturale}: in tutti e tre i casi,
la scelta \textit{on-premise} non è nata da una preferenza tecnologica ma da un vincolo
normativo --- GDPR per KU Leuven e TU Delft, sovranità giurisdizionale per CINECA. Questo
dato conferma la tesi centrale della presente ricerca: per gli atenei europei, il
\textit{compliance by design} coincide frequentemente con l'\textit{on-premise by design}.

Il secondo pattern è la \textbf{condivisione consortile come soluzione al problema del TCO}:
sia CINECA sia i consorzi di ricerca in cui opera KU Leuven distribuiscono i costi fissi
dell'infrastruttura su un numero elevato di istituzioni partecipanti, rendendo sostenibile
un livello di capacità computazionale che nessun ateneo singolo potrebbe permettersi. Questo
pattern è replicabile su scala italiana attraverso i meccanismi già esistenti --- CINECA, rete
GARR --- e su scala europea attraverso l'infrastruttura EuroHPC \parencite{EuroHPC_Leonardo}.

Il terzo pattern è la \textbf{separazione tra il dato grezzo e l'output del modello}: in tutti
e tre i casi, l'architettura garantisce che il dato personale grezzo non abbandoni il perimetro
istituzionale, mentre solo le predizioni aggregate o i parametri del modello (nel caso del
\textit{federated learning}) circolano tra i nodi. Questa separazione è la traduzione
architetturale del principio di \textit{data minimisation} del GDPR e del requisito di
supervisione umana dell'AI Act: l'output del modello è verificabile dall'operatore umano,
mentre il dato grezzo è protetto.

% ============================================================
\section{Implicazioni strategiche e roadmap per gli atenei italiani}
\label{sec:implicazioni-strategiche}

\subsection{Dalla governance del dato alla strategia AI: una sequenza obbligata}

L'analisi dei casi europei e dei \textit{framework} di governance del dato convergono su
una conclusione managerialmente rilevante: l'adozione di sistemi di IA \textit{on-premise}
non è una questione tecnologica che precede quella organizzativa, ma una
\textbf{sequenza obbligata} in cui la maturità organizzativa è il prerequisito della
sostenibilità tecnologica. Un ateneo che investa in hardware GPU senza aver prima classificato
i propri dati, nominato un CDO operativo e certificato l'infrastruttura secondo ISO/IEC\,27001
non sta adottando l'IA \textit{on-premise}: sta acquisendo capacità computazionale non
governata, che il primo audit di compliance renderà inutilizzabile.

\textcite{Torichnyi2025StrategicMgmt} identificano nella strategia di \textit{innovative leadership}
--- caratterizzata dall'integrazione tra visione strategica di vertice, capacità organizzative
di middle management e investimenti tecnologici mirati --- il modello più efficace per la
trasformazione digitale della PA. Applicato agli atenei, questo modello implica che la scelta
\textit{on-premise} sia deliberata dal CDA (livello strategico), tradotta in un piano operativo
dal Delegato alla Transizione Digitale (livello tattico) e implementata da strutture tecniche
con competenze di \textit{data engineering} e \textit{MLOps} (livello operativo). La coerenza
verticale tra questi tre livelli è la condizione per cui la scelta architetturale produce
valore istituzionale e non si riduce a un adempimento burocratico.

\textcite{Pham2025MACSRG} dimostrano, in un'analisi empirica su PA di diversi paesi, che
l'integrazione tra sistemi di controllo della gestione e governance del rischio è positivamente
correlata alla resilienza organizzativa. Per gli atenei che adottano IA \textit{on-premise},
la governance del rischio algoritimco --- identificazione, valutazione e mitigazione dei rischi
connessi ai modelli di IA in produzione --- diventa parte integrante del sistema di controllo
manageriale, non un adempimento normativo separato.

\subsection{Tre fattori abilitanti trasversali}

L'analisi comparativa dei casi europei e dei \textit{framework} di governance identifica tre
fattori abilitanti che compaiono sistematicamente nelle esperienze di maggiore successo.

Il primo è l'\textbf{infrastruttura condivisa sovrana}: la rete GARR per la connettività
e CINECA per il calcolo forniscono agli atenei italiani una base infrastrutturale
\textit{on-premise} a costi consortili. Il modello europeo EuroHPC, di cui Leonardo è
parte, estende questa logica alla scala continentale. La condivisione dell'infrastruttura
riduce il TCO dell'\textit{on-premise} a livelli comparabili con il \textit{cloud} per gli
atenei medi, eliminando uno dei principali argomenti contro la scelta
\textit{on-premise} \parencite{Madhavan2025Storage}.

Il secondo fattore è la \textbf{governance del dato istituzionalizzata}: DPO operativo,
CDO nominato, classificazione dei dati completata e politiche di \textit{data retention}
coerenti con i requisiti dell'AI Act. Senza questi prerequisiti organizzativi, nessuna
infrastruttura tecnologica --- per quanto potente --- è in grado di supportare sistemi di IA
conformi. \textcite{Aguilar2026} dimostrano che negli enti PA europei che hanno adottato
approcci \textit{data-centric} alla governance dell'IA --- cioè centrati sulla qualità e
sul controllo del dato piuttosto che sulla sola potenza computazionale --- i risultati in
termini di affidabilità e accountability dei sistemi sono significativamente superiori.

Il terzo fattore è la \textbf{compliance come leva strategica}: l'AI Act Annex III non
è solo un obbligo normativo per gli atenei che utilizzano IA in ambito educativo; è anche
una \textbf{barriera all'ingresso} per i fornitori che non sono in grado di garantire
la conformità richiesta. Un ateneo che costruisce un sistema di IA \textit{on-premise}
conforme all'AI Act dispone di un vantaggio competitivo nella partecipazione a bandi di
ricerca europei --- dove la protezione dei dati è una condizione di ammissibilità --- e
nella collaborazione con partner istituzionali che richiedono garanzie sulla sovranità
del dato.

\subsection{Raccordo verso il Capitolo\,4}

I pattern architetturali emersi dall'analisi dei casi universitari europei --- \textit{federated
learning}, architettura ibrida con inferenza \textit{on-premise}, modello consortile ---
trovano riscontro, con le necessarie differenziazioni settoriali, nelle esperienze di
altre PA europee analizzate nel Capitolo\,4. La convergenza tra le esigenze di compliance
degli atenei e quelle di altre amministrazioni pubbliche --- ministeri, agenzie regolatorie,
enti locali --- conferma che la scelta \textit{on-premise} risponde a una logica manageriale
e normativa trasversale, che supera i confini del settore accademico e si configura come un
modello replicabile su scala europea.

In questo senso, l'università non è solo un caso studio: è, come anticipato nella sezione
introduttiva, un \textbf{laboratorio istituzionale} dove la tensione tra innovazione
algoritmica, sovranità del dato e compliance normativa si manifesta con una chiarezza analitica
che rende le soluzioni adottate immediatamente generalizzabili ad altri contesti
della PA digitale europea.
